\section{From Dynamic Externality to Modal Materiality}
\subsection{Nominal Modal Consequence}

The E05 modal study selected three coalitions mechanically from E03/E04 before
examining any damping result: A7--A8 as the largest reproducible
curvature-dominated pair, A3--A6 as the largest reproducible
feedback-dominated pair, and A2--A7--A8 as the largest reproducible triple.
Each candidate achieved valid mode tracking at all eight holdout points, with
consistent signs.

The median finite damping externalities were
$-1.84336\times10^{-4}$ for A7--A8,
$-4.09123\times10^{-4}$ for A3--A6, and
$+1.18769\times10^{-5}$ for A2--A7--A8. A7--A8 and A3--A6 also produced
median total joint damping changes of $-0.00391371$ and $-0.00565123$,
respectively. Thus their portfolios change modal damping, but most of that
change is additive at the tracked-mode level. None of the 24 candidate/seed
points reached $|\ext{S}\zeta|\ge0.0025$. The nominal materiality claim was
therefore rejected under its preregistered rule.

This failure is not a contradiction. The characteristic functional integrates
the complete descriptor over a frequency band, while a damping gate observes
one tracked modal family. A large value of $\ext{S}\Phi$ may arise from
algebraic re-equilibration, pole motion in a well-damped family, or a
distributed collection of small shifts. Only the last stage asks whether the
non-additive effect consumes a decision-relevant margin.

\subsection{Load-Conditioned Stress}

E05C tested whether modal composition amplified as load increased toward the
first 0.85-pu voltage boundary. Eight new seeds, all 15 retained coalitions,
and nine stress levels were evaluated with 100\% valid externality and
tracking rates. The feasible endpoint multiplier ranged from 1.05225 to
1.083125 across seeds.

Define the real-part composition ratio
\begin{equation}
\rho_{\mathrm{comp}}=
\frac{\Re\{\lambda_{S}-\lambda_{S}^{(<|S|)}\}}
{\max(|\Re\{\lambda_{S}^{(<|S|)}\}|,\epsilon_\lambda)},
\label{eq:rho}
\end{equation}
where $\lambda_{S}^{(<|S|)}$ is the lower-order inclusion--exclusion
prediction and $\epsilon_\lambda$ is the registered numerical floor. The
largest observed $\rho_{\mathrm{comp}}$ was 0.00211758, and the largest
candidate median at the final stress point was 0.00197909. No point reached
the material damping threshold. There were no hard composition failures and
no 5\% screen crossings attributable to composition before the voltage
boundary. Both load-conditioned amplification and material consequence were
rejected.

Although $\rho_{\mathrm{comp}}$ was correlated with conditioning and
remaining margin, those regressions were explanatory and did not change a
candidate decision. A statistically monotone trend cannot substitute for the
frozen magnitude gate.

\subsection{Weak-Grid Attempt: Baseline Ineligibility}

Conventional short-circuit ratio (SCR) is useful context for GFL operation in
weak networks, but multi-infeed interactions can make single-POI strength
metrics incomplete \cite{xin2016gscr,nerc2017lowstrength}. Converter control
can also introduce weak-grid dynamics that require more than static strength
screening \cite{fan2019type4}. E05D therefore preregistered a grid-strength
path with an immutable eligibility rule: the empty-coalition baseline had to
retain at least 5\% damping in the true minimum-damping 0.1--30 Hz mode.

At $\kappa_{\mathrm{grid}}=1$, all eight new baselines already violated that
screen. Their minimum damping ranged from 0.0416536 to 0.0452751. The
corresponding conventional single-POI values were
SCR36=2.60275--2.64401, SCR37=2.57783--2.65366, and
SCR38=1.32141--1.40561. Because there was no eligible starting point, no safe
stress endpoint existed and no E05D coalition outcome was inspected.

The only valid status is \emph{unresolved--baseline ineligible}. It is neither
a rejection nor support of weak-grid materiality. The 5\% rule was not
changed, the stress direction was not reversed, and no replacement candidate
search was opened.

\subsection{Critical-Mode Baseline Audit}

After E05D closed, a separate descriptive audit asked why the eligibility
screen and the externality search referred to such different margins. It used
only $\kappa_{\mathrm{grid}}=1$, the eight E05D seeds, all 15 frozen E05B
coalitions, and the original A1--A8 definitions. It did not create a new gate.

For each empty baseline, the audit selected the true minimum-damping
positive-imaginary eigenvalue in 0.1--30 Hz, retained its left/right vectors,
condition number, BMAC, and state and descriptor participations, and tracked
that exact mode through every coalition subset. Global one-to-one matching
required BMAC at least 0.90, with registered 5-, 9-, and 17-point homotopies.
All 200 unique subset tracks and all 120 coalition/seed comparisons succeeded.

All eight limiting modes were synchronous-electromechanical. Their frequencies
were 1.21973--1.22906 Hz, damping ratios were 0.0416536--0.0452751, and
synchronous-machine/governor/exciter states carried
0.98619--0.98982 of normalized descriptor participation. By contrast, the 15
externality-dominant development modes were PLL/current-control families with
damping 0.49390--0.99168. Figure~\ref{fig:separation} exposes both the damping
and physical-family separation.

The limiting family is a coherent multi-machine electromechanical mode rather
than a single converter-local mode. Median descriptor endogeneity is led by
the rotor speed/angle channels of GENROU32, GENROU30, GENROU34, GENROU31, and
GENROU33; the median angle-channel shares are 0.1450, 0.0849, 0.0816, 0.0727,
and 0.0547, respectively. This state anatomy supports the family label without
promoting it to a causal intervention target.

\begin{figure*}[t]
\centering
\includegraphics[width=0.94\textwidth]{figure_4_externality_vs_margin_mode.pdf}
\caption{Externality-dominant development modes versus the actual
margin-setting mode at $\kappa_{\mathrm{grid}}=1$. Left: the development
modes are much more strongly damped than the 4.17--4.53\% critical-mode band.
Right: they are converter-control/PLL dominated, whereas the critical mode is
synchronous-electromechanical near 1.22 Hz.}
\label{fig:separation}
\end{figure*}

Across all 120 coalition/seed cases, the largest critical-mode
$|\ext{S}\zeta|$ was $5.14674\times10^{-6}$. The maximum
$\rho_{\mathrm{comp}}$ was $1.39286\times10^{-4}$ for A7--A8 on seed D02,
where $\ext{\{7,8\}}\Re\lambda=4.56767\times10^{-5}\ \mathrm{s}^{-1}$.
The most stabilizing ratio was $-9.63916\times10^{-6}$ for A5--A8, and the
median absolute ratio over all cases was $6.38782\times10^{-6}$. These values
are descriptive evidence for a mode-family mismatch, not a post hoc
reinterpretation of E05D.

Table~\ref{tab:materiality} summarizes the resulting boundary. Structural
externalities survive every exactness and reproducibility check, while
operational materiality does not survive the registered consequence tests.

\begin{table}[t]
\caption{Externality Existence Versus Modal Materiality}
\label{tab:materiality}
\centering
\footnotesize
\begin{tabular}{p{0.35\columnwidth}p{0.24\columnwidth}p{0.22\columnwidth}}
\toprule
Evidence & Value & Interpretation\\
\midrule
A7--A8 spectral externality & $-0.00411$ median holdout & Exists; algebraic dominated\\
A3--A6 pole externality & $-3.36\times10^{-4}$ & Exists; pole dominated\\
Nominal $|\ext{S}\zeta|$ & $4.09\times10^{-4}$ median & Below 0.0025 gate\\
Load-stress $\rho_{\mathrm{comp}}$ & 0.00211758 maximum & No material point\\
Weak-grid campaign & 8/8 baselines below 5\% & Ineligible; unresolved\\
Critical-mode audit & $5.15\times10^{-6}$ max $|\ext{S}\zeta|$ & Below registered gate on limiting family\\
\bottomrule
\end{tabular}
\end{table}
